\section{Introduction}

El reconocimiento facial es un problema clásico de la visión computacional,
pues trae consigo aplicaciones en seguridad, autenticación y análisis biométrico.
En la actualidad, el avance de las redes neuronales profundas
ha permitido lograr tasas de reconocimiento altas que antes eran inimaginables
\cite{taigman2014deepface,schroff2015facenet}; sin embargo, esto lleva consigo una
limitación estructural fundamental: su funcionamiento interno es ininterpretable
\cite{rudin2019stop}, lo que dificulta en gran medida su análisis, corrección y
aplicación, más aún en contextos en los cuales es necesaria la transparencia.
En sistemas de identificación judicial, si el sistema de reconocimiento facial es usado
para identificar a una persona como sospechosa, resulta difícil impugnar una decisión
errónea con un sistema basado en redes neuronales, pues no existe una estructura
interpretable que auditar ni una forma directa de demostrar cómo y por qué el sistema
falló, lo que obliga a depender de pruebas empíricas exhaustivas. Por ello, se ha
propuesto preferir modelos inherentemente interpretables en decisiones de alto riesgo,
como las de justicia penal y salud \cite{rudin2019stop}.
En sistemas de infraestructura crítica, como el control de plantas de energía,
instalaciones militares o bancos, es deseable poder auditarlos de manera formal y no
solo a base de experimentos. En el ámbito de la investigación, para construir sobre un
resultado es necesario saber cómo y por qué funciona; un modelo ininterpretable no
permite derivar mejoras teóricas.

El presente trabajo propone un sistema alternativo basado en álgebra lineal: el
algoritmo Eigenfaces, planteado por Turk y Pentland en 1991 \cite{turk1991eigenfaces}.
Fue una de las soluciones pioneras computacionalmente viables al problema del
reconocimiento facial: dada la limitada capacidad de cómputo de la época, reemplazó la
comparación píxel por píxel por la proyección de los rostros en un subespacio de menor
dimensión.
A diferencia de los sistemas que hacen uso de redes neuronales, los cuales se denominan
``caja negra'' debido a su ininterpretabilidad, Eigenfaces exhibe los patrones de
variación utilizados para el reconocimiento de rostros, permitiendo así el análisis e
interpretación del proceso de reconocimiento. El sistema es construido en Python haciendo
uso exclusivo de la biblioteca NumPy con el fin de conservar la mayor transparencia
posible de la estructura algebraica.
Hoy en día, pese a ser en general menos preciso que los sistemas basados en redes
neuronales, Eigenfaces conserva su importancia como modelo de estudio: al exponer de
manera explícita los fundamentos matemáticos del reconocimiento de patrones (reducción
dimensional, análisis de varianza y proyección en subespacios), permite construir
intuición sobre el problema antes de abordar sistemas más complejos y comprender los
conceptos en profundidad, en lugar de limitarse a aplicar lo ya establecido.

Se establecen los siguientes objetivos específicos para el presente trabajo:
(1) implementación del algoritmo Eigenfaces haciendo uso exclusivo de operaciones
de álgebra lineal sobre NumPy;
(2) análisis del efecto del número de componentes $k$ sobre la tasa de reconocimiento
del sistema; y
(3) evaluación de la viabilidad de un modelo matemáticamente interpretable mediante
su tasa de reconocimiento y costo computacional, frente a líneas base sencillas.
Los objetivos en cuestión buscan contrastar si la interpretabilidad matemática y la
eficiencia computacional pueden coexistir con un reconocimiento aceptable, y si un
sistema desarrollado a partir de fundamentos matemáticos sólidos puede ser una
herramienta válida tanto para su aplicación como para la formación y comprensión de
los principios del reconocimiento facial.

El presente sistema se limita al reconocimiento facial en condiciones controladas,
haciendo uso de la base de datos AT\&T Database of Faces, la cual contiene 400 imágenes
de 40 personas, tomadas en posición frontal con variaciones de iluminación y
expresión; por tanto, no se abordará el reconocimiento facial en tiempo real ni la
detección facial previa, y se asume que las imágenes de entrada ya están recortadas y
alineadas.
Además, el sistema se enmarca en el aprendizaje automático clásico, específicamente en
el reconocimiento de patrones basado en métodos algebraicos y estadísticos. Se espera
que, a diferencia de otros enfoques modernos que suelen requerir grandes volúmenes de
datos y recursos computacionales para aprender sus representaciones, Eigenfaces opere
con conjuntos de datos pequeños; esta expectativa se contrasta experimentalmente en la
metodología. De confirmarse, sería una alternativa viable en escenarios con recursos
limitados o donde la recopilación masiva de datos biométricos presenta restricciones
éticas o logísticas.
Para finalizar, el presente trabajo no busca competir con otros modelos en precisión,
sino demostrar que es posible construir un sistema de reconocimiento facial viable a
partir de herramientas matemáticas elementales. El valor del trabajo consiste en la
transparencia del proceso: la justificación algebraica es explícita, cada resultado
puede analizarse y cada limitación es predecible a partir de la estructura del sistema.